\section{Outlook}\label{ss:Outlook}
In this review, we put emphasis on the roles of quantum inequalities of entanglement entropy in constraining the dynamics of RG flows in QFTs.
In particular, the monotonicity of $\CC$-functions built from entanglement entropy follows from the strong subadditivity, which should be contrasted with Zamolodchikov's $c$-theorem and the $a$-theorem in two- and four-dimensions whose proofs are based on unitarity in QFT.
Clearly the strong subadditivity assumes almost an equal role with unitarity when applied to QFTs, and indeed it is a consequence of the Hermicity of density matrices in finite-dimensional quantum mechanical systems \cite{Lieb:1973cp,araki1976relative}.
It is thus important to see whether the strong subadditivity can be proved by a conventional method within the canonical framework of unitary and Lorentz invariant QFTs.

Some attempts for $\CC$-theorems in higher dimensions were discussed in Sec. \ref{ss:C_theorem_Higher} where a certain class of monotonic functions were constructed out of entanglement entropy.
They result in the known $\CC$-theorems in less than four-dimensions, but do not lead to the conjectured $F$- and $a$-theorems in higher dimensions.
It deserves further investigation to see if the conjectures can be proved by a similar method to the lower-dimensional case.

It would also be interesting to consider $\CC$-theorems for RG flows across dimensions.
This possibility is anticipated on general grounds, but has attracted less attention so far \cite{Gukov:2015qea,Bobev:2017uzs}.
It is not clear at first sight what can be a measure of degrees of freedom, i.e.\, a $\CC$-function for RG flows across dimensions.
A naive speculation is to use as a measure the generalized $F$ coefficient defined by Eq.\,\eqref{Generalized_F} possibly with a dimension-dependent normalization factor.
It would be intriguing to explore this direction and find an appropriate measure that may or may not have a relation to entanglement entropy.